\documentclass[11pt,a4paper]{article}

\usepackage[utf8]{inputenc}
\usepackage[english]{babel}
\usepackage{fontspec}
\setmainfont{Times New Roman}
\setsansfont{Arial}
\setmonofont{Consolas}

\usepackage[a4paper,top=1.5cm,bottom=1.6cm,left=1.7cm,right=1.7cm]{geometry}
\usepackage{amsmath,amssymb}
\usepackage{booktabs}
\usepackage{tabularx}
\usepackage{array}
\usepackage{listings}
\usepackage{xcolor}
\usepackage{fancyhdr}
\usepackage{multicol}
\usepackage{enumitem}
\usepackage{tcolorbox}
\usepackage{lastpage}

% Colors for syntax highlighting and frames
\definecolor{codegray}{rgb}{0.96,0.96,0.96}
\definecolor{keywordblue}{rgb}{0.0,0.2,0.6}
\definecolor{stringgreen}{rgb}{0.0,0.5,0.2}
\definecolor{commentgray}{rgb}{0.4,0.4,0.4}
\definecolor{boxborder}{rgb}{0.7,0.7,0.7}

\lstset{
  language=Python,
  backgroundcolor=\color{codegray},
  basicstyle=\ttfamily\small,
  keywordstyle=\color{keywordblue}\bfseries,
  stringstyle=\color{stringgreen},
  commentstyle=\color{commentgray}\itshape,
  showstringspaces=false,
  breaklines=true,
  frame=single,
  rulecolor=\color{boxborder},
  aboveskip=2.5pt,
  belowskip=2.5pt,
  numbers=none,
  columns=fullflexible
}

% Header & Footer
\pagestyle{fancy}
\fancyhf{}
\setlength{\headheight}{14pt}
\lhead{\small \textbf{Course:} Data Analysis with Python (DSAI1005)}
\rhead{\small \textbf{Exam (Weeks 1 -- 4) --- 60 Minutes}}
\lfoot{\small Faculty of Data Science \& AI --- National Economics University}
\rfoot{\small Page \thepage/6}
\renewcommand{\headrulewidth}{0.5pt}
\renewcommand{\footrulewidth}{0.5pt}

\newcommand{\answerbox}[2]{
  \vspace{1mm}
  \begin{tcolorbox}[
    colback=white,
    colframe=gray!60,
    boxrule=0.6pt,
    arc=2pt,
    height=#1,
    left=3mm, right=3mm, top=2mm, bottom=2mm
  ]
  \small\textit{\color{gray}#2}
  \end{tcolorbox}
  \vspace{1mm}
}

\begin{document}

% =========================================================================
% PAGE 1: HEADER + SCORE TABLE + ANSWER GRID + QUESTIONS 1 TO 4
% =========================================================================
\noindent
\begin{minipage}[t]{0.54\textwidth}
  \centering
  \textbf{NATIONAL ECONOMICS UNIVERSITY}\\
  \textbf{FACULTY OF DATA SCIENCE \& AI}\\
  \rule{4.5cm}{0.4pt}\\[1.5mm]
  \textbf{MIDTERM EXAMINATION (PAPER-BASED)}\\
  \textbf{Course:} Data Analysis with Python\\
  \textbf{Scope:} Core Lectures 1 -- 4 (Intro, NumPy, Pandas, Viz)\\
  \textbf{Course Code:} DSAI1005 --- \textbf{Duration:} 60 minutes\\
  \textit{\small (No personal laptops, textbooks, or notes allowed)}
\end{minipage}
\hfill
\begin{minipage}[t]{0.43\textwidth}
  \begin{tabular}{|p{2.2cm}|p{2.2cm}|p{2.2cm}|}
    \hline
    \centering \textbf{Section I} & \centering \textbf{Section II} & \centering \textbf{TOTAL} \tabularnewline
    \hline
    & & \\[0.8cm]
    \hline
    \multicolumn{3}{|l|}{\small \textit{Examiner 1:}} \tabularnewline
    \multicolumn{3}{|l|}{\small \textit{Examiner 2:}} \tabularnewline
    \hline
  \end{tabular}
\end{minipage}

\vspace{2mm}
\noindent
\textbf{Full Name:} \dotfill{} \textbf{Student ID:} \dotfill{} \textbf{Exam Code:} \textbf{201-EN}\\
\textbf{Class / Cohort:} \dotfill{} \textbf{Exam Room:} \dotfill{} \textbf{Seat No:} \dotfill

\vspace{2mm}
\noindent
\textbf{MULTIPLE-CHOICE ANSWER SHEET} \textit{(Write letters A, B, C, or D in the corresponding boxes)}:
\begin{center}
\begin{tabular}{|c|c|c|c|c|c|c|c|c|c|c|}
  \hline
  \textbf{Question} & \textbf{1} & \textbf{2} & \textbf{3} & \textbf{4} & \textbf{5} & \textbf{6} & \textbf{7} & \textbf{8} & \textbf{9} & \textbf{10} \\
  \hline
  \textbf{Answer} & & & & & & & & & & \\[3mm]
  \hline
  \hline
  \textbf{Question} & \textbf{11} & \textbf{12} & \textbf{13} & \textbf{14} & \textbf{15} & \textbf{16} & \textbf{17} & \textbf{18} & \textbf{19} & \textbf{20} \\
  \hline
  \textbf{Answer} & & & & & & & & & & \\[3mm]
  \hline
\end{tabular}
\end{center}

\vspace{-3mm}
\hrule
\vspace{1.5mm}

\section*{SECTION I: MULTIPLE-CHOICE QUESTIONS (5.0 pts --- 0.25 pts each)}

\begin{enumerate}[leftmargin=*,itemsep=2pt]

\item What is the core objective of the ``Exploratory Data Analysis (EDA)'' stage within the standardized 6-step data analysis lifecycle?
\begin{enumerate}[label=\Alph*.]
  \item Immediately train complex black-box machine learning models to maximize benchmark accuracy.
  \item Understand data structure, test hypotheses, uncover underlying patterns, and detect anomalies/outliers.
  \item Optimize the physical storage engine and indexes in a relational database server.
  \item Export raw tabular files to third-party vector illustration software for visual touch-ups.
\end{enumerate}

\item If every observation in a numerical dataset is multiplied by a constant factor of 3, the standard deviation of the new dataset will:
\begin{multicols}{2}
\begin{enumerate}[label=\Alph*.]
  \item Increase by a factor of 9
  \item Remain strictly unchanged
  \item Increase by a factor of 3
  \item Equal the square root of the original value
\end{enumerate}
\end{multicols}

\item In a standard Tukey Box Plot, an observation $x$ is formally classified as a potential outlier if it falls outside which of the following intervals (where $\text{IQR} = Q_3 - Q_1$)?
\begin{multicols}{2}
\begin{enumerate}[label=\Alph*.]
  \item $[Q_1 - 1.5\times\text{IQR}, \; Q_3 + 1.5\times\text{IQR}]$
  \item $[Q_1 - 3.0\times\text{IQR}, \; Q_3 + 3.0\times\text{IQR}]$
  \item $[\text{Mean} - 1.5\times\text{Std}, \; \text{Mean} + 1.5\times\text{Std}]$
  \item $[0, \; Q_3 + \text{IQR}]$
\end{enumerate}
\end{multicols}

\item The Pearson correlation coefficient between two continuous variables $X$ and $Y$ evaluates to $r = 0$. Which of the following statements is most statistically sound?
\begin{enumerate}[label=\Alph*.]
  \item $X$ and $Y$ are completely independent and share no mathematical relationship whatsoever.
  \item There is no linear association between $X$ and $Y$, although a strong non-linear relationship may still exist.
  \item Whenever $X$ increases by 1 unit, $Y$ remains fixed at 0.
  \item Both variables $X$ and $Y$ must follow a standard normal distribution.
\end{enumerate}

\end{enumerate}

\newpage
% =========================================================================
% PAGE 2: QUESTIONS 5 TO 12 (8 QUESTIONS)
% =========================================================================
\begin{enumerate}[leftmargin=*,itemsep=2.5pt,resume]

\item What array is generated by the NumPy expression \texttt{np.arange(3, 15, 4)}?
\begin{multicols}{4}
\begin{enumerate}[label=\Alph*.]
  \item \texttt{[3, 7, 11, 15]}
  \item \texttt{[3, 7, 11]}
  \item \texttt{[4, 8, 12]}
  \item \texttt{[3, 6, 9, 12]}
\end{enumerate}
\end{multicols}

\item Consider the following NumPy snippet:
\begin{lstlisting}
import numpy as np
m = np.array([[1, 2, 3], [4, 5, 6], [7, 8, 9]])
print(m[::-1, 1])
\end{lstlisting}
What output is printed to the console?
\begin{multicols}{4}
\begin{enumerate}[label=\Alph*.]
  \item \texttt{[2, 5, 8]}
  \item \texttt{[8, 5, 2]}
  \item \texttt{[9, 6, 3]}
  \item \texttt{[7, 4, 1]}
\end{enumerate}
\end{multicols}

\item Array $A$ has shape \texttt{(5, 1, 3)} and array $B$ has shape \texttt{(4, 1)}. Under NumPy broadcasting rules, what is the shape of the resulting array from \texttt{A * B}?
\begin{multicols}{4}
\begin{enumerate}[label=\Alph*.]
  \item \texttt{(5, 4, 3)}
  \item \texttt{(5, 1, 3)}
  \item \texttt{(4, 5, 3)}
  \item Raises \texttt{ValueError}
\end{enumerate}
\end{multicols}

\item Which NumPy function performs element-wise conditional substitution, returning elements chosen from $X$ or $Y$ depending on a boolean condition (equivalent to a vectorized ternary operator)?
\begin{multicols}{4}
\begin{enumerate}[label=\Alph*.]
  \item \texttt{np.select()}
  \item \texttt{np.where()}
  \item \texttt{np.filter()}
  \item \texttt{np.choose()}
\end{enumerate}
\end{multicols}

\item For a two-dimensional array \texttt{arr} with shape \texttt{(4, 5)}, what is the shape of the result returned by \texttt{np.sum(arr, axis=0)}?
\begin{multicols}{4}
\begin{enumerate}[label=\Alph*.]
  \item \texttt{(4,)}
  \item \texttt{(5,)}
  \item \texttt{(1,)}
  \item \texttt{(4, 5)}
\end{enumerate}
\end{multicols}

\item Which primary Pandas data structure represents a one-dimensional labeled array capable of holding any data type?
\begin{multicols}{4}
\begin{enumerate}[label=\Alph*.]
  \item \texttt{DataFrame}
  \item \texttt{Series}
  \item \texttt{IndexList}
  \item \texttt{NDArray}
\end{enumerate}
\end{multicols}

\item Given the Pandas Series \texttt{s = pd.Series([10, 20, 30, 40], index=['a', 'b', 'c', 'd'])}, what elements are returned by the slice \texttt{s.loc['a':'c']}?
\begin{multicols}{2}
\begin{enumerate}[label=\Alph*.]
  \item Only elements with index \texttt{'a'} and \texttt{'b'}
  \item All three elements with index \texttt{'a'}, \texttt{'b'}, and \texttt{'c'}
  \item A syntax error because string labels cannot be sliced
  \item Only the single element with index \texttt{'c'}
\end{enumerate}
\end{multicols}

\item In Pandas, which method is specifically designed to impute or replace missing entries (\texttt{NaN}) with a designated value or summary statistic (e.g., column mean)?
\begin{multicols}{4}
\begin{enumerate}[label=\Alph*.]
  \item \texttt{df.dropna()}
  \item \texttt{df.fillna()}
  \item \texttt{df.replace\_na()}
  \item \texttt{df.impute()}
\end{enumerate}
\end{multicols}

\end{enumerate}

\newpage
% =========================================================================
% PAGE 3: QUESTIONS 13 TO 20 (8 QUESTIONS)
% =========================================================================
\begin{enumerate}[leftmargin=*,itemsep=2.5pt,resume]

\item Which statement correctly filters records in DataFrame \texttt{df} where \texttt{credit\_score >= 700} OR \texttt{income > 50000000}?
\begin{enumerate}[label=\Alph*.]
  \item \texttt{df[(df['credit\_score'] >= 700) or (df['income'] > 50000000)]}
  \item \texttt{df[(df['credit\_score'] >= 700) | (df['income'] > 50000000)]}
  \item \texttt{df[(df['credit\_score'] >= 700) || (df['income'] > 50000000)]}
  \item \texttt{df.filter(credit\_score >= 700 or income > 50000000)}
\end{enumerate}

\item Given the following Pandas snippet:
\begin{lstlisting}
import pandas as pd
df = pd.DataFrame({'Dept': ['Sales', 'IT', 'IT', 'Sales'], 'Salary': [30, 50, 70, 40]})
print(df.groupby('Dept')['Salary'].agg(['min', 'max']).loc['IT', 'max'])
\end{lstlisting}
What value is printed?
\begin{multicols}{4}
\begin{enumerate}[label=\Alph*.]
  \item \texttt{30}
  \item \texttt{50}
  \item \texttt{70}
  \item \texttt{60}
\end{enumerate}
\end{multicols}

\item When executing an inner merge between two DataFrames using \texttt{pd.merge(df1, df2, on='id', how='inner')}, a row appears in the output DataFrame if and only if:
\begin{enumerate}[label=\Alph*.]
  \item The key \texttt{'id'} is present in either \texttt{df1} or \texttt{df2}.
  \item The key \texttt{'id'} matches simultaneously in both \texttt{df1} and \texttt{df2}.
  \item \texttt{df1} contains strictly more rows than \texttt{df2}.
  \item The key column contains strictly unique non-null integers.
\end{enumerate}

\item How many distinct subplot axes (\texttt{Axes}) are created by the call \texttt{fig, axes = plt.subplots(2, 3)} in Matplotlib?
\begin{multicols}{4}
\begin{enumerate}[label=\Alph*.]
  \item 5 axes
  \item 6 axes
  \item 2 axes
  \item 3 axes
\end{enumerate}
\end{multicols}

\item Which Seaborn function is best suited to estimate and visualize a smooth, continuous probability density curve (Kernel Density Estimation) for a financial variable?
\begin{multicols}{4}
\begin{enumerate}[label=\Alph*.]
  \item \texttt{sns.kdeplot()}
  \item \texttt{sns.scatterplot()}
  \item \texttt{sns.countplot()}
  \item \texttt{sns.barplot()}
\end{enumerate}
\end{multicols}

\item In Seaborn, which high-level function automatically renders an $n \times n$ grid showing pairwise scatter plots across all quantitative features and univariate distributions along the diagonal?
\begin{multicols}{4}
\begin{enumerate}[label=\Alph*.]
  \item \texttt{sns.jointplot()}
  \item \texttt{sns.pairplot()}
  \item \texttt{sns.heatmap()}
  \item \texttt{sns.catplot()}
\end{enumerate}
\end{multicols}

\item In Matplotlib, which utility function is commonly called prior to display or export to automatically adjust subplot margins, preventing axes labels and titles from overlapping?
\begin{multicols}{4}
\begin{enumerate}[label=\Alph*.]
  \item \texttt{plt.align\_axis()}
  \item \texttt{plt.tight\_layout()}
  \item \texttt{plt.clean\_view()}
  \item \texttt{plt.auto\_padding()}
\end{enumerate}
\end{multicols}

\item To compare the mean values of a numerical metric across categorical groups while displaying confidence intervals (error bars) by default, which Seaborn function is optimal?
\begin{multicols}{4}
\begin{enumerate}[label=\Alph*.]
  \item \texttt{sns.lineplot()}
  \item \texttt{sns.barplot()}
  \item \texttt{sns.histplot()}
  \item \texttt{sns.heatmap()}
\end{enumerate}
\end{multicols}

\end{enumerate}

\newpage
% =========================================================================
% PAGE 4: SECTION II - QUESTION 1 (NUMPY) + LARGE ANSWER BOX
% =========================================================================
\section*{SECTION II: SHORT-ANSWER \& CODE COMPREHENSION (5.0 pts)}

\subsection*{Question 1: Vectorized Computations \& Multidimensional Arrays with NumPy (1.0 pt) --- [Week 2]}
Consider the exam score matrix of 4 students across 3 subjects (Math, Probability, Programming) and the corresponding course credit weights:
\begin{lstlisting}
import numpy as np

scores = np.array([
    [8.0, 7.0, 9.0],
    [6.0, 8.0, 7.0],
    [9.0, 9.0, 8.5],
    [5.0, 6.0, 6.5]
])  # Shape: (4, 3)

weights = np.array([0.2, 0.3, 0.5])  # Shape: (3,)
\end{lstlisting}

\noindent
\textbf{Tasks:}
\begin{enumerate}[label=\textbf{\alph*)}]
  \item \textbf{(0.5 pts)} Name the fundamental NumPy mechanism that enables direct multiplication between \texttt{scores} of shape \texttt{(4, 3)} and \texttt{weights} of shape \texttt{(3,)}. Write the exact numerical matrix for \texttt{weighted\_scores = scores * weights}.
  \item \textbf{(0.5 pts)} Write a single NumPy statement to calculate the 1D array \texttt{gpa} (representing the weighted grade point average for each student by summing along the row axis of \texttt{weighted\_scores}), and determine the exact output printed by: \quad \texttt{print(scores[gpa >= 8.0, 0])}
\end{enumerate}

\answerbox{12.0cm}{
\textbf{Work Area --- Question 1:}\\[3mm]
\textbf{a) Name of mechanism:} \dotfill\\[4mm]
\textbf{Resulting matrix weighted\_scores:}\\[28mm]
\textbf{b) NumPy statement to compute gpa:}\\
gpa = \dotfill\\[5mm]
\textbf{Exact output printed by the print statement:}\\
\dotfill\\[3mm]
\textit{(Brief reasoning if applicable):} \dotfill
}

\newpage
% =========================================================================
% PAGE 5: QUESTION 2 (PANDAS) + LARGE ANSWER BOX
% =========================================================================
\subsection*{Question 2: Data Cleaning \& Aggregation with Pandas (1.5 pts) --- [Week 3]}
You are given a Pandas DataFrame \texttt{df\_trans} representing retail banking transaction logs with columns:
\texttt{'branch'} (branch code), \texttt{'amount'} (transaction value in VND), \texttt{'fee'} (service charge), and \texttt{'status'} (transaction outcome: \texttt{'SUCCESS'} or \texttt{'FAILED'}). The column \texttt{'fee'} contains several missing entries (\texttt{NaN}).

\noindent
\textbf{Tasks (write correct and concise Pandas code):}
\begin{enumerate}[label=\textbf{\alph*)}]
  \item \textbf{(0.5 pts)} Write Pandas code to filter rows satisfying two conditions: \texttt{'status'} equals \texttt{'SUCCESS'} and \texttt{'amount'} $\ge 5\,000\,000$ VND. Simultaneously, impute all missing values (\texttt{NaN}) in the \texttt{'fee'} column of this filtered subset with $0$.
  \item \textbf{(1.0 pts)} Using the filtered dataset from part (a), construct a chained Pandas expression (using \texttt{groupby} and \texttt{agg}) to group by branch (\texttt{'branch'}), compute: total transaction volume named \texttt{'total\_amount'}, average fee named \texttt{'avg\_fee'}, and transaction count named \texttt{'trans\_count'}. Sort the final summary table in descending order of \texttt{'total\_amount'}.
\end{enumerate}

\answerbox{14.0cm}{
\textbf{Work Area --- Question 2:}\\[3mm]
\textbf{a) Filtering rows and imputing missing values (NaN):}\\
df\_filtered = \dotfill\\[4mm]
df\_filtered['fee'] = \dotfill\\[7mm]
\textbf{b) Groupby, aggregation, and sorting expression:}\\
df\_branch\_summary = (\dotfill\\[5mm]
\hspace*{3.5cm}\dotfill\\[5mm]
\hspace*{3.5cm}\dotfill\\[5mm]
\hspace*{3.5cm}\dotfill\\[5mm]
\hspace*{3.5cm}\dotfill)\\[5mm]
\dotfill
}

\newpage
% =========================================================================
% PAGE 6: QUESTION 3 (MATPLOTLIB) + QUESTION 4 (EDA & VIZ) + END
% =========================================================================
\subsection*{Question 3: Multi-plot Layouts with Matplotlib (1.0 pt) --- [Week 4]}
Given a DataFrame \texttt{df\_kpi} tracking monthly performance with columns \texttt{'month'} (integers 1 to 12) and \texttt{'profit'} (net profit in million VND).

\noindent
\textbf{Tasks (complete the code template by filling in the blanks):}
\begin{enumerate}[label=\textbf{\alph*)}]
  \item \textbf{(0.5 pts)} Write the statement to initialize a figure of size $12 \times 4$ inches containing 1 row and 2 columns of subplots, storing the objects in \texttt{fig} and array \texttt{axes}.
  \item \textbf{(0.5 pts)} On \texttt{axes[0]}, render a line plot (\texttt{plot}) illustrating monthly profit trajectory (blue line with circular markers \texttt{'o'}). On \texttt{axes[1]}, render a bar chart (\texttt{bar}) of monthly profit. Finally, apply layout optimization.
\end{enumerate}

\answerbox{3.8cm}{
\textbf{Work Area --- Question 3:}\\
import matplotlib.pyplot as plt\\[1.5mm]
fig, axes = \dotfill\\[2.5mm]
axes[0].plot(\dotfill)\\[2.5mm]
axes[1].bar(\dotfill)\\[2.5mm]
\dotfill \quad \textit{(\# Command to prevent overlapping labels/titles)}
}

\subsection*{Question 4: Exploratory Visualization \& Business Insights (EDA) (1.5 pts) --- [Weeks 1 \& 4]}
A bank's Risk Analytics division evaluates transaction data to detect card fraud. The DataFrame \texttt{df\_cards} contains: \texttt{'tx\_amount'} (monetary value), \texttt{'card\_type'} (categorical: \texttt{'Credit'} or \texttt{'Debit'}), and fraud flag \texttt{'is\_fraud'} ($1$ for fraudulent, $0$ for legitimate).

\noindent
\textbf{Tasks:}
\begin{enumerate}[label=\textbf{\alph*)}]
  \item \textbf{(0.75 pts)} Choose the most appropriate statistical chart type to inspect the distribution and detect potential extreme outliers of \texttt{'tx\_amount'} across fraud labels, stratified by card type. Write a single line of \textbf{Seaborn} (\texttt{sns}) code to render this visualization.
  \item \textbf{(0.75 pts)} The resulting plot reveals: For \texttt{'Debit'} cards, both legitimate and fraudulent transactions cluster at low values with no significant outliers. Conversely, for \texttt{'Credit'} cards, fraudulent transactions display an extreme outlier spread from 50 to 200 million VND, whereas legitimate ones remain under 10 million VND.  
  \textit{Formulate \textbf{2 concise analytical takeaways / policy recommendations (under 50 words)} to guide the bank's automated fraud prevention rule.}
\end{enumerate}

\answerbox{5.0cm}{
\textbf{Work Area --- Question 4:}\\[1.5mm]
\textbf{a) Chosen chart type:} \dotfill\\[2mm]
\textbf{Seaborn statement:} \dotfill\\[3mm]
\textbf{b) Analytical Takeaways \& Policy Recommendations:}\\
- Takeaway 1: \dotfill\\[2mm]
\hspace*{2.2cm}\dotfill\\[2mm]
- Policy 2: \dotfill\\[2mm]
\hspace*{2.2cm}\dotfill
}

\vspace{1mm}
\begin{center}
\textbf{------------------- END OF EXAM -------------------}\\
\textit{\small (Proctors will not answer questions regarding exam content. Hand in all exam sheets upon completion.)}
\end{center}

\end{document}
